\subsection{Recovery after five missing neighbour counts}
\label{sec:five-erasure240}
The four-count result can be extended by a finite calculation involving
five missing positions. The calculation concerns arbitrary binary vertex
selections; it does not require the selected vertices to form a snake.

\begin{lemma}\label{lem:five-matrix240}
Let $p_0,\ldots,p_4$ be distinct vertices of one parity in $Q_9$,
and suppose that some pair has Hamming distance eight. Define $M$ by
\[
M_{ii}=128,\qquad
M_{ij}=\begin{cases}
-16,&\Ham(p_i,p_j)=2,\\
8,&\Ham(p_i,p_j)=4,\\
-10,&\Ham(p_i,p_j)=6,\\
35,&\Ham(p_i,p_j)=8
\end{cases}\quad(i\ne j).
\]
If $w\in\mathbb Z^5$, $s\in\{-1,0,1\}^5$, and $Mw=315s$, then $w=0$.
\end{lemma}
\begin{proof}
We give the geometric reduction and the finite arithmetic certificate.
Translate one endpoint of a distance-eight pair to zero and permute
coordinates so that the other is $(1,1,1,1,1,1,1,1,0)$. The remaining
three positions have even weight. In the first eight coordinates,
record their three bits as one of the eight possible columns. Permuting
these coordinates preserves the fixed pair, so their multiplicities
give all required representatives. The ninth bit of each position is
determined by even parity. Consequently the $\binom{15}{7}=6435$
compositions of eight into eight parts cover every configuration in
the lemma. Removing repeated positions leaves 5160 descriptions.
Relabelling the five positions gives 372 distinct distance matrices.
Only the distance matrix is needed to determine $M$; this reduction
does not assert that every equal-distance configuration is cube-isometric.

For each of these 372 matrices, the saved calculation gives an integer
adjugate $B$ and a positive determinant $D$ satisfying
\[
BM=DI.
\]
Thus $Dw_i=315(Bs)_i$. Put $q=D/\gcd(D,315)$. Since the two remaining
factors are coprime after division by the gcd, integrality requires
\[
q\mid(Bs)_i\quad(0\le i\le4).
\]
In particular, if $q>\sum_j|B_{ij}|$, the ternary restriction on $s$
forces $(Bs)_i=0$, and hence $w_i=0$.

The finite matrices separate as follows:
\begin{center}
\begin{tabular}{rrl}
\toprule
Matrices & Rows with $q>\sum_j|B_{ij}|$ & Remaining argument\\
\midrule
334 & 5 & All coordinates of $w$ vanish.\\
5 & 4 & The last diagonal equation forces zero.\\
33 & 0 & Check the 243 ternary right-hand sides.\\
\bottomrule
\end{tabular}
\end{center}
For the middle row, four unknowns already vanish. The remaining row is
$128w_i=315s_i$. Parity and $s_i\in\{-1,0,1\}$ force $s_i=0$, and then
$w_i=0$.

For each of the last 33 matrices, direct integer evaluation of the
five divisibility conditions for all $3^5=243$ vectors $s$ leaves only
$s=0$. The adjugate identity then gives $w=0$. The saved certificate
lists every matrix, determinant, and adjugate, so these are explicit
finite statements rather than an assumption about an optimizer.
The complete native calculation also checks all
$372\cdot243=90396$ matrix/right-hand-side combinations, with no bound
imposed on the integer unknowns $w$~\cite{research240}.
\end{proof}

\begin{proposition}\label{prop:five-erasure240}
The neighbour counts at any 507 of the 512 vertices of $Q_9$ determine
a binary vertex selection uniquely. If the selected parity side is
known, any 251 of the 256 opposite-side counts suffice.
\end{proposition}
\begin{proof}
Subtract the indicators of two selections agreeing on the observed
counts. Their difference $f$ takes values in $\{-1,0,1\}$, and $Af$
vanishes outside at most five positions. At most four positions are
handled by Proposition~\ref{prop:four-erasure239}.

If the five positions have mixed parity, restrict $f$ to each parity
side. Adjacency sends each restriction to the opposite side, where
its count differences are supported on at most four positions.
The four-count result makes both restrictions zero.

It remains to consider five positions of one parity. If they contain
no distance-eight pair, Proposition~\ref{prop:eight-erasure239} applies.
Otherwise evaluate the cube inverse at their antipodes. With
$w_i=Af(p_i)$ and $s_i=f(\overline{p_i})$, its coefficients give exactly
$Mw=315s$ from Lemma~\ref{lem:five-matrix240}. Therefore $Af=0$, and the
inverse gives $f=0$. For a selection on a known parity side, the 256
counts on that side are automatically zero, leaving only the five
unobserved opposite-side counts.
\end{proof}

The arithmetic and geometric verification have different scopes.
The native audit covers the 6435 column descriptions and the 90396
exact arithmetic cases. An independent enumeration of all
$\binom{254}{3}=2699004$ triples of actual even vertices outside the
fixed pair gives exactly the same 372 distance matrices. A separate source comparison confirms that
the generated 372 Lean statements have the matrices and coefficients
of this catalogue. Each generated proof checks its adjugate identities
and finite divisibility implication, then concludes that every integer
solution is zero. All 372 implications have passed Lean checking, with their
source-bound proof receipts retained. The catalogue coverage, its connection to arbitrary cube
positions, and the mixed-parity reduction remain written arguments
supported by native computation. The fully assembled Lean theorem on
the actual cube still covers three missing counts.

Consequently two different binary selections must change at least six
neighbour counts. Changing one selected vertex changes exactly nine
counts, so the minimum lies between six and nine. A possible failure
with six through eight missing counts must contain a distance-eight
pair, by Proposition~\ref{prop:eight-erasure239}. These are constraints
on reconstruction of vertex selections. They produce neither a longer
snake nor an improved global snake-length bound.
